\chapter{诊断、维护和恢复}

\section{doctor}

常规诊断：

\begin{lstlisting}
ccb doctor
\end{lstlisting}

生成 bundle：

\begin{lstlisting}
ccb doctor --output
ccb doctor --output /tmp/ccb-doctor
\end{lstlisting}

\cmd{doctor --bundle} 已不支持，使用 \cmd{doctor --output}。

\section{runtime 诊断}

\begin{lstlisting}
ccb doctor ps
ccb doctor --runtime
ccb ps
\end{lstlisting}

用于查看 backend、agent runtime、health 和相关状态。

\section{日志诊断}

\begin{lstlisting}
ccb doctor logs
ccb doctor --logs
ccb logs <agent>
\end{lstlisting}

日志适合排查 provider 启动、pane、runtime 和 completion 问题。

\section{存储诊断}

\begin{lstlisting}
ccb doctor storage
ccb doctor storage --json
\end{lstlisting}

当 queue/inbox 显示 summary missing 或 summary error 时，优先运行 storage diagnostics。

\section{maintenance}

\begin{lstlisting}
ccb maintenance
ccb maintenance status
ccb maintenance tick
ccb maintenance tick --force
ccb maintenance tick --no-dispatch
ccb maintenance schedule --after 5m --reason "follow up"
\end{lstlisting}

maintenance heartbeat 可用于周期性检查 agent health。默认 assessor 是 \cfg{ccb_self}，默认 policy 是 \cfg{report_only}。

\cmd{ccb maintenance schedule} 不是只读查看命令；它必须带 \cmd{--after <duration>}，并会写入下一次 heartbeat follow-up。当前 \cmd{ccb maintenance enable} 和 \cmd{ccb maintenance disable} 只返回 config-authority 提示；要启停 heartbeat，请编辑 \cfg{[maintenance.heartbeat].enabled}。这里的 config-authority 说明不改变 \src{.ccb/ccb.config} 必须使用 \cfg{version = 2} 的规则。

\section{reload}

修改配置后先验证：

\begin{lstlisting}
ccb config validate
\end{lstlisting}

预览 reload：

\begin{lstlisting}
ccb reload --dry-run
\end{lstlisting}

执行 reload：

\begin{lstlisting}
ccb reload
\end{lstlisting}

\section{restart}

重启单个 agent：

\begin{lstlisting}
ccb restart archi
\end{lstlisting}

\cmd{restart all} 不支持。restart 只接受单个 agent，并由 ccbd busy gate 保护；如果 agent 正在执行、排队、投递回复或等待 callback continuation，会被拒绝。需要全项目处理时使用 reload、kill 或重新启动 CCB。

\section{kill}

普通停止：

\begin{lstlisting}
ccb kill
\end{lstlisting}

强制：

\begin{lstlisting}
ccb kill --force
ccb kill -f
\end{lstlisting}

不要手动删除项目 \src{.ccb/agents/*} 或 \src{.ccb/ccbd/*} 来代替 kill。

\section{cleanup}

\begin{lstlisting}
ccb cleanup
ccb cleanup --legacy-provider-caches
\end{lstlisting}

用于在当前 ccbd 已停止且没有待处理 ask 时清理可重建缓存和运行残留。普通命令只处理当前项目；增加
\cmd{--legacy-provider-caches} 后，还会删除 manifest 校验通过、且记录的项目根目录已经不存在的旧版
Claude/Gemini 项目级缓存。该选项不会删除仍存在项目的缓存，也不会删除 provider session、auth、mailbox
或 runtime authority。cleanup 不应被当作常规启动步骤。

新版本安装成功后，\cmd{ccb update} 会自动执行同一安全边界的迁移：
已删除项目中逐 Provider manifest 校验通过的缓存会清理；活跃项目不被
停止，其缓存会推迟到该项目下一次成功 \cmd{ccb kill}。普通
\cmd{ccb kill} 也会清理当前项目的旧缓存。无法验证的路径保持不动，
\cmd{ccb update --no-cache-cleanup} 可跳过单次更新迁移。

清理前可运行 \cmd{ccb doctor storage}。输出会列出旧项目级 Provider 缓存的路径、存在状态和字节数，
并列出新的用户级 Provider 缓存路径和存在状态。用户级缓存由多个项目共享，常规项目诊断不会递归统计
它的大小；这些外部缓存不混入项目 \src{.ccb} 的 \cmd{storage_total_bytes}。

\section{fault injection}

高级测试：

\begin{lstlisting}
ccb fault list
ccb fault arm archi --task-id task1 --reason api_error --count 1 --error "drill"
ccb fault clear all
\end{lstlisting}

fault injection 适合测试失败路径，不适合普通用户日常使用。
